\subsection{ETL Pipeline Implementation}

The ingestion layer is a standalone Python script (\texttt{db/etl\_script.py}) that consumes the CSV bundle generated by Synthea and writes it into a normalised PostgreSQL schema accessed by the rest of the application through SQLModel. The script does not depend on the FastAPI process: it is executed once per dataset refresh and produces the relational tables that all later layers (clustering, pattern mining, statistics, SQL agent) read from.

Database connectivity is centralised in \texttt{db/db\_connector.py}, which exposes a single SQLAlchemy \texttt{engine}, a \texttt{SessionLocal} factory, a \texttt{get\_db()} FastAPI dependency, and a thin \texttt{\_db\_retrieve(query)} helper that returns raw queries as \texttt{list[dict]}. Credentials are read from environment variables (\texttt{POSTGRES\_USERNAME}, \texttt{POSTGRES\_PASSWORD}, \texttt{POSTGRES\_HOST}, \texttt{POSTGRES\_PORT}, \texttt{POSTGRES\_DB}) loaded via \texttt{python-dotenv}, so the same code base runs against local and remote databases without modification.

Each Synthea CSV is read with \texttt{pandas.read\_csv} in a streaming-friendly manner, lower-cased on the column axis, and validated against a target SQLModel definition before being passed to \texttt{DataFrame.to\_sql} with \texttt{if\_exists="append"} and \texttt{method="multi"} to batch inserts. Tables that are conceptually many-to-many (e.g.\ patient/provider) are populated by a second pass over the \texttt{encounter} table that produces the \texttt{patient\_provider} junction with \texttt{DISTINCT} on the relevant pair of identifiers.

Computed columns described in Section~5.1.1 are materialised through pure SQL after the raw load, using \texttt{ALTER TABLE \dots ADD COLUMN} followed by an \texttt{UPDATE \dots FROM} pattern. For instance, \texttt{patient.age} is derived as:

\begin{verbatim}
UPDATE patient
SET age = DATE_PART('year', AGE(COALESCE(deathdate, CURRENT_DATE), birthdate));
\end{verbatim}

\noindent and counts such as \texttt{provider.encounters} are filled with a correlated aggregate update so that downstream analytical queries can avoid the join entirely. Foreign keys are added at the end of the script with \texttt{ALTER TABLE \dots ADD CONSTRAINT}, after the data has been loaded, to keep the bulk insert phase fast.

At application startup, \texttt{app/main.py} runs a small set of idempotent migrations that add the columns required for the analytical features (\texttt{ai\_analysis JSONB} on \texttt{cluster\_meta}, \texttt{cluster}, and \texttt{prefix\_span\_pattern}, plus \texttt{pca\_image BYTEA} on \texttt{cluster\_meta}). The statements use \texttt{ADD COLUMN IF NOT EXISTS} so a redeployment over an existing database is a no-op:

\begin{verbatim}
with engine.begin() as conn:
    conn.execute(text(
        "ALTER TABLE cluster_meta ADD COLUMN IF NOT EXISTS ai_analysis JSONB"
    ))
\end{verbatim}

The principal libraries involved in the ETL layer are \texttt{pandas} for CSV parsing and dataframe manipulation, \texttt{SQLAlchemy}/\texttt{SQLModel} for schema declaration and connection pooling, and \texttt{psycopg2-binary} as the PostgreSQL driver. No ORM-level inserts are used for the bulk load: all heavy writes go through \texttt{to\_sql} or raw \texttt{text()} statements, which keeps the import of a full Synthea dataset (\(\approx\)\,20{,}000 patients, several million rows across all tables) within a few minutes on a developer machine.
